\documentclass[11pt]{article}
\usepackage[margin=1in]{geometry}
\usepackage{amsmath,amssymb,amsthm,mathtools}
\usepackage[colorlinks=true,allcolors=blue]{hyperref}
\setlength{\emergencystretch}{2em}
\newtheorem{theorem}{Theorem}
\newtheorem{proposition}[theorem]{Proposition}
\newcommand{\R}{\mathbb R}
\newcommand{\C}{\mathbb C}
\newcommand{\N}{\mathbb N}
\newcommand{\J}{\mathcal J}
\newcommand{\CC}{\mathcal C}
\newcommand{\Dom}{\operatorname{Dom}}
\title{Independent native-mass derivation of the arithmetic Jacobi commutant}
\author{}
\date{20 September 2026}
\begin{document}
\maketitle

\section{Unchanged data and the proved envelope}

Retain the actual original packet divisor
\[
 \rho=\tfrac12+\delta+i\gamma,\quad 0<\delta<\tfrac12,\quad\gamma>2,
 \quad h(s)=\prod_{z\in\{\rho,\bar\rho,1-\rho,1-\bar\rho\}}(s-z)^m,
\]
with its full multiplicity $m\geq1$, $g=2\xi$, and
\[
 w_h(y)=\frac{|(g/h)(1/2+iy)|^2}{2\pi},\qquad
 m_k=w_h^{*k},\qquad \mu_h=\int_\R w_h(y)\,dy.
\]
Here $k=4l+1\geq9$ is the original tensor order. These are the
original programme objects, and no existence assertion for an
off-line packet is added. The roots of $h$ are closed under conjugation
with unchanged multiplicities, so $h$ has real coefficients. The
original completed zeta function satisfies
$g(\bar s)=\overline{g(s)}$. Thus $w_h(-y)=w_h(y)$; convolution
preserves this equality and $m_k$ is even. Put
\[
 c=k/2,\quad q=[1+k(m-1)](k+1)^2,\qquad
 \chi(S)=\prod_{a,b=0}^k
 [S-c-(2a-k)\delta-i(2b-k)\gamma]^{1+k(m-1)}.
\]
The full finite quotient will use $N\geq q-1$.

The reference measure and its mass are exactly
\[
 d\sigma(y)=\rho_\sigma(y)\,dy
 =\frac{|\Gamma(1/4+iy/2)|^2}{2\pi}\,dy,
 \qquad M_\sigma=\sqrt{2\pi}.
\]
In particular, this measure is not replaced by a probability measure.
The already proved CAI16--18 envelope supplies
\begin{equation}\label{eq:envelope}
 a_k(1+y^2)^{-M}\rho_\sigma(y)\leq m_k(y)
                       \leq A_k\rho_\sigma(y),\qquad M=21+4m,
\end{equation}
where the retained constants are
\[
 \alpha=\pi/2,\quad p=8m-9/2,\quad B_h=42+8m,\quad
 \vartheta_h=\int_{-1}^1w_h(y)\,dy,
 \quad M_\alpha=\int_\R e^{\alpha y}w_h(y)\,dy,
\]
\[
 c_\Gamma=\sqrt2e^{-7/6},\qquad C_\Gamma=\sqrt{2\sqrt5}e^{2/3},
\]
\[
 a_k=\frac{c_h\vartheta_h^{k-3}e^{-\alpha(k-3)}(k-2)^{-B_h}}
              {C_\Gamma2^{B_h/2}},\qquad
 A_k=\frac{C_h^{\rm tilt}}{c_\Gamma}k^{p+1}M_\alpha^{k-1}.
\]
The constants $c_h,C_h^{\rm tilt}$ are the original CAI lower and
upper envelope constants. Their proof is not replaced by a hypothesis:
the source is the retained programme provider
\texttt{CAI\_\allowbreak COMPLETE\_\allowbreak PREREQUISITE.tex},
equations CAI16--18.
The all-degree consequence, with
$\Delta_N=[1+4(N+M)^2]^M$, is
\begin{equation}\label{eq:finite-envelope}
 \frac{a_k}{\Delta_N}\|P\|_\sigma^2
       \leq\int|P(y)|^2m_k(y)\,dy\leq A_k\|P\|_\sigma^2
       \qquad(\deg P\leq N).
\end{equation}
The notation $\Delta_N$ here distinguishes the retained scalar CAI
constant from the source-coordinate matrix $D_N$ below.

The human literature input is Franti\v{s}ek \v{S}tampach and Pavel
\v{S}\v{t}ov\'{i}\v{c}ek, \emph{Spectral representation of some
weighted Hankel matrices and orthogonal polynomials from the Askey
scheme}; see
\href{https://arxiv.org/abs/1811.05820v1}{arXiv:1811.05820v1}.
We use Theorem \texttt{thm:commut\_\allowbreak matr\_\allowbreak Jac},
its following remark and proof,
and Theorem \texttt{thm:def\_H}. The following is a direct reconstruction
for the actual multiplier, including the mass absent from that
source's probability-measure convention.

\section{Gamma polynomials and the native Jacobi operator}

Define the monic polynomials by
\begin{equation}\label{eq:generating}
 G(y,t)=\sum_{n\geq0}p_n(y)\frac{t^n}{n!}
       =(1+t^2)^{-1/4}e^{y\arctan t}.
\end{equation}
Every branch here is its analytic branch at $t=0$. To verify the
mass and norms directly, the beta integral gives, for $a>0$,
\[
 \frac{\Gamma(a+i\tau)\Gamma(a-i\tau)}{\Gamma(2a)}
 =\int_\R\frac{e^{i\tau v}}{(2\cosh(v/2))^{2a}}\,dv.
\]
This follows from $u=e^v$ in the beta integral on $(0,\infty)$.
Fourier inversion, followed by analytic continuation in the strip
where the Gamma envelope is integrable, gives at $a=1/4$ and $y=2\tau$
\begin{equation}\label{eq:mgf-sigma}
 \int_\R e^{zy}\,d\sigma(y)
 =\sqrt{2\pi}(\cos z)^{-1/2},\qquad |\Re z|<\pi/2.
\end{equation}
In detail, inversion gives the Fourier transform at real frequency
$v$ as $2\Gamma(1/2)/(2\cosh v)^{1/2}$; setting $v=-iz$ yields
$\sqrt{2\pi}(\cos z)^{-1/2}$. The CAI17 Gamma envelope bounds the
integrand on every closed narrower strip and justifies holomorphy
and this continuation. At zero this proves the stated mass.

For sufficiently small $t,u$, the cosine addition formula and
\eqref{eq:mgf-sigma} imply
\[
 \int_\R G(y,t)G(y,u)\,d\sigma(y)
 =M_\sigma(1-tu)^{-1/2}
 =M_\sigma\sum_{n\geq0}\frac{(1/2)_n}{n!}t^nu^n.
\]
All coefficient extractions may be taken on fixed small complex
circles, dominated by $Ce^{\varepsilon|y|}$ with $\varepsilon<\pi/2$.
Consequently
\[
 \int p_m p_n\,d\sigma=\delta_{mn}\gamma_n,
 \qquad\gamma_n=M_\sigma n!(1/2)_n.
\]
Define, without altering the measure,
\[
 \varphi_n=\frac{p_n}{\sqrt{\gamma_n}},\qquad
 r_n=\sqrt{M_\sigma}\varphi_n
     =\frac{p_n}{\sqrt{n!(1/2)_n}}.
\]
Thus $\varphi_0=M_\sigma^{-1/2}$, $r_0=1$, and
$\int r_mr_n\,d\sigma=M_\sigma\delta_{mn}$.
The $r_n$ are not orthonormal in $d\sigma$.

Differentiating \eqref{eq:generating} gives
$(1+t^2)\partial_tG=(y-t/2)G$ and hence
\[
 p_{n+1}=yp_n-n(n-1/2)p_{n-1},\qquad p_0=1,\quad p_1=y.
\]
Therefore
\begin{equation}\label{eq:recurrence}
 y\varphi_n=b_n\varphi_{n+1}+b_{n-1}\varphi_{n-1},\qquad
 b_n=\sqrt{(n+1)(n+1/2)},\quad b_{-1}=0,
\end{equation}
and the same recurrence holds for $r_n$. Also $p_n(-y)=(-1)^np_n(y)$.
The original source's Meixner--Pollaczek generating function gives
the exact dictionary
\[
 p_n(y)=n!P_n^{(1/4)}(y/2;\pi/2),\qquad
 r_n(y)=\sqrt{\frac{n!}{(1/2)_n}}P_n^{(1/4)}(y/2;\pi/2).
\]
Both the coordinate $y/2$ and every normalization factor remain.

The polynomials are complete in $L^2(\sigma)$. Indeed, if $f$ is
orthogonal to all polynomials, then $f\,d\sigma$ is a finite complex
measure, and Cauchy--Schwarz together with
\eqref{eq:mgf-sigma} shows that
$\int f(y)e^{zy}\,d\sigma(y)$ is holomorphic for $|\Re z|<\pi/4$.
Every derivative at zero is zero, so analytic uniqueness makes
the transform zero throughout that connected strip. Fourier
uniqueness on the imaginary axis makes $f\,d\sigma=0$, hence $f=0$.
The identical argument with the finite measure $(1+y^2)d\sigma$
proves polynomial density in that weighted space as well.

It follows that
\[
 U_\sigma:\ell^2(\N_0)\longrightarrow L^2(\sigma),\qquad
 U_\sigma e_n=\varphi_n
\]
is unitary. The operator $J=U_\sigma^*M_yU_\sigma$ is self-adjoint,
and its matrix $\J$ has the entries \eqref{eq:recurrence}, with zero
diagonal. Polynomial density for $(1+y^2)d\sigma$ means that the
polynomials are a core of $M_y$ in its graph norm. Thus the finite
vectors are a core of $J$, so this is the unique self-adjoint closure
of the stated Jacobi matrix on finite vectors.

\section{The original bounded multiplier and exact first column}

Set $R_k=m_k/\rho_\sigma$. By \eqref{eq:envelope},
\[
 0<a_k(1+y^2)^{-M}\leq R_k(y)\leq A_k.
\]
Therefore multiplication by this exact density ratio is bounded,
positive and self-adjoint on $L^2(\sigma)$, and
\begin{equation}\label{eq:operator}
 \CC=U_\sigma^*M_{R_k}U_\sigma,\qquad
 0\leq\CC\leq A_k I,\qquad
 \CC_{mn}=\int_\R\varphi_m(y)\varphi_n(y)m_k(y)\,dy.
\end{equation}
This is not merely a formally defined matrix. The bounded multiplier
preserves $\Dom M_y$, since
$\|yR_k f\|_2\leq A_k\|yf\|_2$. Consequently $\CC$ preserves
$\Dom J$ and $J\CC=\CC J$ there. Entrywise the same statement is
\[
 b_{m-1}\CC_{m-1,n}+b_m\CC_{m+1,n}
 =b_{n-1}\CC_{m,n-1}+b_n\CC_{m,n+1}.
\]
Each side is a finite sum; it also follows at once by integrating
\eqref{eq:recurrence} against the other polynomial and $m_k$.

Put $\alpha_j=\CC_{j0}$. The exact native-mass expression is
\begin{equation}\label{eq:alpha-moment}
 \alpha_j=\frac1{M_\sigma}\int r_j(y)m_k(y)\,dy
 =\frac{\int p_j(y)m_k(y)\,dy}
             {M_\sigma\sqrt{j!(1/2)_j}}.
\end{equation}
In particular $\alpha_0=\mu_h^k/M_\sigma$ by Tonelli.

\begin{theorem}\label{thm:first}
For every $m,n\geq0$ the complete entry is the finite sum
\begin{equation}\label{eq:first}
 \CC_{mn}=\sum_{j=|m-n|}^{m+n}\alpha_j[r_j(\J)]_{mn}.
\end{equation}
Here $r_j(\J)$ is evaluated in the full Jacobi matrix, not in a
premature finite compression.
\end{theorem}
\begin{proof}
The recurrence and $r_0=1$ give $r_n(\J)e_0=e_n$ by induction.
Indeed, for $n=0$ the recurrence gives $r_1(\J)e_0=e_1$; if the
assertion holds at $n,n-1$, then
$b_n r_{n+1}(\J)e_0=\J e_n-b_{n-1}e_{n-1}=b_ne_{n+1}$.
Expand the degree-$m+n$ polynomial $r_mr_n$ in the $r_j$ basis.
Orthogonality shows that its coefficient at $r_j$ is
\[
 d_{mnj}=\frac1{M_\sigma}\int r_mr_nr_j\,d\sigma
       =\int\varphi_m r_j\varphi_n\,d\sigma
       =[r_j(\J)]_{mn}.
\]
The last equality follows either from the multiplication operator
or by applying \eqref{eq:recurrence} finitely many times. Its
integral is zero when $j>m+n$ by orthogonality. It is also zero
when $m+j<n$ or $n+j<m$, by applying orthogonality to $r_n$ or
$r_m$ respectively. Thus only $|m-n|\leq j\leq m+n$ occurs.
Integrate this polynomial identity against $m_k/M_\sigma$ and
use \eqref{eq:alpha-moment}. The result is \eqref{eq:first}.
\end{proof}

Evenness of the unchanged weight gives $\alpha_j=0$ for odd $j$ and
$\CC_{mn}=0$ for odd $m+n$. The first column also reconstructs the
actual multiplier in $L^2(\sigma)$, with the exact mass factor:
\begin{equation}\label{eq:mass-caveat}
 U_\sigma\CC e_0=\frac{R_k}{\sqrt{M_\sigma}},\qquad
 R_k=\sqrt{M_\sigma}U_\sigma\CC e_0
       =\sum_{j\geq0}\alpha_jr_j\quad\text{in }L^2(\sigma).
\end{equation}
To prove the last convergence, expand the vector $\CC e_0\in\ell^2$
in its coordinates $\alpha_j$ and apply the unitary $U_\sigma$.
In particular the multiplier itself is not $U_\sigma\CC e_0$ in
the original non-probability convention. No pointwise or
operator-norm convergence of the displayed series is asserted.

Define $M_h(z)=\int e^{zy}w_h(y)\,dy$. CAI16 gives
$M_\alpha<\infty$. Since the weight is even,
$\int e^{\alpha|y|}w_h\leq2M_\alpha$. Tonelli and convolution give
\[
 \int e^{zy}m_k(y)\,dy=M_h(z)^k,
 \qquad |\Re z|<\alpha.
\]
For real $z$ this is the absolutely convergent product integral;
the same domination gives the complex version and holomorphy.
For $|t|\leq1/4$,
$|\arctan t|\leq\operatorname{arctanh}(1/4)<\alpha$.
We can therefore integrate \eqref{eq:generating} on any smaller
complex circle and extract coefficients. This proves exactly
\begin{equation}\label{eq:alpha-generating}
 \boxed{\alpha_j=
 \frac{j!}{M_\sigma\sqrt{j!(1/2)_j}}
 [t^j]\left((1+t^2)^{-1/4}M_h(\arctan t)^k\right).}
\end{equation}
The factor $1/M_\sigma=1/\sqrt{2\pi}$ is necessary.

\section{The finite-cutoff sign and its full boundary vector}

Let $P_N$ project to $e_0,\ldots,e_N$, and write
\[
 C_N=P_N\CC P_N,\quad J_N=P_N\J P_N,\quad
 c_N=(\CC_{0,N+1},\ldots,\CC_{N,N+1})^T.
\]
All are now finite coordinate matrices/vectors. Compressing the
full entrywise commutation identity gives
\[
 0=[J_N,C_N]+P_N\J(I-P_N)\CC P_N
                    -P_N\CC(I-P_N)\J P_N.
\]
The unique crossing Jacobi edge is $b_N$. Since $\CC$ is real
symmetric, the two extra terms are respectively $b_Ne_Nc_N^*$
and $b_Nc_Ne_N^*$. Hence
\begin{equation}\label{eq:cutoff}
 \boxed{[J_N,C_N]=b_N(c_Ne_N^*-e_Nc_N^*).}
\end{equation}
This proves the stated sign. Formula \eqref{eq:first} computes
$C_N$ using the coefficients through index $2N$. It computes $c_N$
using those through index $2N+1$. Because the native odd coefficients
vanish, $\alpha_{2N+1}$ is zero; retaining it in the index range is
correct and harmless.

In fact evenness gives $(c_N)_N=\CC_{N,N+1}=0$, so $c_N\perp e_N$.
If $c_N=0$, the finite matrices commute. If $c_N\ne0$, the
commutator vanishes on the orthogonal complement of
$\operatorname{span}\{e_N,c_N\}$ and acts on the ordered orthonormal
basis $e_N,c_N/\|c_N\|$ there as
\[
 \begin{pmatrix}0&-b_N\|c_N\|\\b_N\|c_N\|&0\end{pmatrix}.
\]
Thus its rank is two and its norm is exactly $b_N\|c_N\|$.
These are identities, not new estimates on the vector $c_N$.

The distinction between $P_Nr_j(\J)P_N$ and $r_j(J_N)$ is essential.
For example $r_2(y)=(y^2-1/2)/\sqrt{3/2}$, and at $N=1$,
\[
 [r_2(\J)]_{11}=\sqrt6,\qquad [r_2(J_1)]_{11}=0.
\]
Indeed $(\J^2)_{11}=b_0^2+b_1^2=1/2+3$, whereas
$(J_1^2)_{11}=b_0^2=1/2$. The missing term is the path from $1$
to $2$ and back to $1$. Accordingly
$\CC_{11}=\alpha_0+\sqrt6\alpha_2$, not merely $\alpha_0$.

\section{Exact original source and observation coordinates}

For $F(S)=\sum_{j=0}^N a_jS^j$, the monomial coefficient vector
of the same function $F(c+iy)$ is $Ta$, where
\[
 T_{rj}=\binom jr c^{j-r}i^r\ (r\leq j),\qquad T_{rj}=0\ (r>j).
\]
This is the binomial theorem with no conjugation or mass change.
Let $B$ have the monomial coefficient vector of $\varphi_n$ as
column $n$. It is upper triangular with diagonal
$1/\sqrt{\gamma_n}$, while $T$ has diagonal $i^n$.
Both matrices are invertible. If $x$ denotes the Gamma-frame
coefficient vector of $F(c+iy)$, equality of the two polynomial
expansions is $Bx=Ta$. Hence
\begin{equation}\label{eq:source}
 x=D_Na,\qquad D_N=B^{-1}T,\qquad
 (D_N)_{nn}=i^n\sqrt{\gamma_n}.
\end{equation}
Integration in the unchanged arithmetic measure gives
\begin{equation}\label{eq:gram}
 H_{S,N}^{\rm ar}=D_N^*C_ND_N.
\end{equation}
For every $x\in\C^{N+1}$, \eqref{eq:finite-envelope} applied to
$\sum x_n\varphi_n$ gives the concrete positive bounds
\begin{equation}\label{eq:CN-bounds}
 \frac{a_k}{\Delta_N}I\leq C_N\leq A_kI.
\end{equation}
Thus the inverse in the next formula exists for every fixed $N$.
No positive lower bound independent of $N$ for the infinite
operator is inferred from the decaying lower density envelope.

Retain the actual observation map
$\mathcal A_S=\Lambda\mathcal J_{S,N}$, including its full physical
Taylor unit, in its original source and observation coordinates.
Its codomain is the original observation image with its retained
Hermitian structure, so it is onto. Its matrix in the Gamma-frame
source coordinates must be
\begin{equation}\label{eq:E}
 \mathcal E_N=\mathcal A_SD_N^{-1},
\end{equation}
since $\mathcal A_Sa=\mathcal A_SD_N^{-1}x$. The symbol
$\mathcal J_{S,N}$ denotes the original class map, not the Jacobi
matrix $J_N$.

\begin{theorem}\label{thm:quotient}
For every observed vector $u$ the original attained minimum is
\begin{equation}\label{eq:quotient}
 \min_{\mathcal A_Sa=u}a^*H_{S,N}^{\rm ar}a
  =u^*Q_{B,N}u,\qquad
 Q_{B,N}=(\mathcal E_NC_N^{-1}\mathcal E_N^*)^{-1},
\end{equation}
and its unique minimizing source vector is
\begin{equation}\label{eq:lift}
 a_{\min}=D_N^{-1}C_N^{-1}\mathcal E_N^*Q_{B,N}u.
\end{equation}
\end{theorem}
\begin{proof}
Surjectivity of $\mathcal E_N$ makes its adjoint injective, so
$K=\mathcal E_NC_N^{-1}\mathcal E_N^*$ is strictly positive on the
observation space. Put $x_0=C_N^{-1}\mathcal E_N^*K^{-1}u$.
Then $\mathcal E_Nx_0=u$ and
$x_0^*C_Nx_0=u^*K^{-1}u$. Every other lift is $x_0+z$ with
$\mathcal E_Nz=0$, and
$z^*C_Nx_0=(\mathcal E_Nz)^*K^{-1}u=0$. Therefore its energy is
$u^*K^{-1}u+z^*C_Nz$, strictly larger unless $z=0$.
The bijection $x=D_Na$ and \eqref{eq:gram} prove both formulas in
the original source coordinates. Equivalently,
\[
 H_{S,N}^{{\rm ar},-1}=D_N^{-1}C_N^{-1}(D_N^*)^{-1},\qquad
 \mathcal A_SH_{S,N}^{{\rm ar},-1}\mathcal A_S^*
                   =\mathcal E_NC_N^{-1}\mathcal E_N^*.
\]
Thus this is exactly the original quotient metric, not an auxiliary
minimum for a changed source.
\end{proof}

For completeness, \eqref{eq:CN-bounds} also gives the exact matrix
comparison
\[
 \frac{a_k}{\Delta_N}(\mathcal E_N\mathcal E_N^*)^{-1}
       \leq Q_{B,N}\leq
 A_k(\mathcal E_N\mathcal E_N^*)^{-1}.
\]
This follows by inverting the two positive bounds for $C_N$,
congruence with $\mathcal E_N$, and inversion once again. It does
not replace $\mathcal E_N$ by another observation map.

\section{Explicit finite checks and retained phase problem}

Write $\nu_j=\int y^jm_k(y)\,dy$ and
$\eta_j=\int y^jw_h(y)\,dy$, so $\eta_0=\mu_h$.
The first polynomial values from the exact recurrence are
\[
 p_0=1,\quad p_1=y,\quad p_2=y^2-\tfrac12,\quad
 p_3=y^3-\tfrac72y,\quad p_4=y^4-11y^2+\tfrac{15}4.
\]
Consequently the first even column entries are
\[
 \alpha_0=\frac{\nu_0}{M_\sigma},\qquad
 \alpha_2=\frac{\nu_2-\nu_0/2}{M_\sigma\sqrt{3/2}},\qquad
 \alpha_4=\frac{\nu_4-11\nu_2+15\nu_0/4}
                   {M_\sigma\sqrt{315/2}}.
\]
Multinomial expansion of the actual convolution, using evenness,
gives
\[
 \nu_0=\mu_h^k,\quad \nu_2=k\eta_2\mu_h^{k-1},\quad
 \nu_4=k\eta_4\mu_h^{k-1}
              +3k(k-1)\eta_2^2\mu_h^{k-2}.
\]
The $3k(k-1)$ factor counts the two distinct factors contributing
second powers: $\binom{k}{2}\binom42=3k(k-1)$.

At cutoff $N=1$, evenness gives
\[
 C_1=\frac1{M_\sigma}\begin{pmatrix}\nu_0&0\\0&2\nu_2\end{pmatrix},
 \quad J_1=\begin{pmatrix}0&1/\sqrt2\\1/\sqrt2&0\end{pmatrix},
 \quad c_1=\begin{pmatrix}\alpha_2\\0\end{pmatrix},\quad b_1=\sqrt3.
\]
Its $(0,1)$ commutator entry is
$(2\nu_2-\nu_0)/(\sqrt2M_\sigma)=\sqrt3\alpha_2$,
confirming both the sign and the mass factor in \eqref{eq:cutoff}.
The degree-two source change is explicitly
\[
 D_2=\sqrt{M_\sigma}
 \begin{pmatrix}
 1&c&c^2-1/2\\
 0&i/\sqrt2&\sqrt2ic\\
 0&0&-\sqrt{3/2}
 \end{pmatrix}.
\]
In particular the source constant $F=1$ has coefficient
$x_0=\sqrt{M_\sigma}$ and energy
$M_\sigma(C_0)_{00}=\mu_h^k$, retaining the original arithmetic mass.

The companion exact symbolic script \texttt{check\_symbolic.py}
checks \eqref{eq:first} for all $0\leq m,n\leq4$ (25 entries),
checks \eqref{eq:cutoff} at $N=0,1,2,3$, checks the generating
polynomials through degree six, and checks
\eqref{eq:source}--\eqref{eq:gram} through degree two. All moment
variables are independent real symbols; the checks even allow
nonzero odd moments. The parameter representing $M_\sigma$ remains
a symbolic positive mass rather than being set to one. These
finite checks supplement, and do not replace, the general proofs.

For the original classes $a_1=[p_N^{\rm ar}]_\chi$,
$a_2=[p_{N+1}^{\rm ar}]_\chi$ and original terminal class $v$, the
retained quantities are therefore exactly
\[
 B_i=(\Lambda a_i)^*Q_{B,N}\Lambda v,\qquad
 \Phi_{B,N}=\frac{2\Re(\overline{B_1}B_2)}{\omega_{N,k}^{\rm ar}}.
\]
The formulas above give a finite representation of the same
quantities and the full Jacobi compression defect. They do not
evaluate or estimate these complex signed products. The physical
Taylor unit, the observation map, and both original classes remain
present; parity of $C_N$ does not remove them.

\end{document}
